\section{A surface-anchored wetting ruler}\label{sec:wetting}
The film engine can also read a coupling-selected thickness. Here the
volume convention requires particular care at height ties. We use
open cubes with $L=n!$ for the exact all-coupling statement. Away
from height ties the same inner gap identity holds along all
sufficiently large open cubes. The anchors and closing modes are
on faces and an edge; this is not an intrinsic bulk-gap realization.

\subsection{Local geometry and energy}
Use the six face flags of Section~\ref{sec:film} and a wetting track
with symbols $F,A,S,D,V,W,O$, denoting floor, signal surroundings,
signal, divider, ceiling and vacuum. Require $F$ exactly at $Z_0$.
Allowed $x$ pairs are
\[
 FF,\ WW,\ OO,\ AA,\ SA,\ DS,\ DD,\ AV,\ SV,\ VS,\ VD.
\]
Allowed $z$ pairs, written in order of decreasing $z$, are
\[
 WS,\ WA,\ WV,\ SF,\ DF,\ VF,\ AA,\ AS,\ SD,\ DD,\ VV,\ OW,\ OO.
\]
Allowed $y$ pairs are $FF,WW,OO$ and every pair within
$\{A,S,D,V\}$. Require the far $Z_1$ face to carry $O$ or $W$.
Every other pair or onsite pattern costs one. The coupling contributes
$g$ per interior symbol $A,S,D,V$, and a $VF$ bond costs $r=1/16$.
All interactions are rational and affine in $g$.

In a clean configuration the floor is $z=0$, a flat ceiling is at
$z=h+1$, $1\leq h\leq L-2$, and vacuum fills the space above it.
The signal cycle cuts each $(x,z)$ interior slice into periods
$h+1$, with a free horizontal phase in each $y$ plane. The
floor itself is not cut. The energy of one such plane, after
subtracting the flag energy, is
\begin{equation}\label{eq:wetting-row}
 E_L(h)=Lgh+r\left\lfloor\frac{L}{h+1}\right\rfloor
\end{equation}
for an optimal phase; an extra divider costs another $r$.
The limiting plane energy per width is
\begin{equation}\label{eq:wetting-phi}
 \phi_g(h)=gh+\frac{r}{h+1},\qquad
 t_h=\frac{r}{(h+1)(h+2)},\qquad
 \phi_g(h)-\phi_g(h+1)=t_h-g.
\end{equation}
Thus off ties $h_*(g)=h$ on $(t_h,t_{h-1})$ and
$h_*(g)=(r/g)^{1/2}+O(1)$.

\begin{lemma}[Wetting selection with commensurate open cubes]
\label{lem:wetting-selection}
Fix sufficiently small positive rational $g$. Along sufficiently
large $L=n!$, the classical minimizing heights are exactly the
minimizers of $\phi_g$. There is one such height off ties, and
two adjacent heights at a tie. All other classical configurations
have a positive separation independent of $L$. The same assertion
off ties holds along all sufficiently large $L$, with a possibly
smaller positive separation than $1/2$.
\end{lemma}
\begin{proof}
The flag identity puts every wrong flag at unit cost. Decompose a
fixed $y$ plane into maximal clean runs of columns with equal floor
and ceiling heights and no violated intervening $x$ bond. A run of
$m$ columns has interior cost $mgh$ and at least
$\lfloor m/(h+1)\rfloor$ divider corners. Its energy is therefore
at least $m\phi_g(h)-r$. A bad column or broken horizontal bond
terminates a run. For $0<g\leq1/8$, a bad column has unit error
cost and nonnegative interior and corner energies, whereas the
minimum plane energy per width is at most $\phi_g(1)\leq1/4$.
Charging each extra run to its terminating error gives
\[
 E(R_y)\geq L\phi_*-r+
             (3/4-r)(\#\text{bad columns}+\#\text{broken }x\text{ bonds}).
\]
Additional $y$ violations only raise the energy. Correct flags and
clean rules are attainable at an optimal height and phase, with
energy at most $L\phi_*$ per plane. A defective plane consequently
costs at least $5/8$ relative to that comparison. Flag errors outside
the plane decomposition are also nonnegative and charged once.

At fixed nontie $g$, put
$\delta(g)=\min_{h\ne h_*}(\phi_g(h)-\phi_*)>0$. The minimum
is effective because $\phi_g(h)\to\infty$. A clean wrong-height
plane lies at least $L\delta(g)-r$ above an optimal one, so it is
excluded for large $L$. A nonoptimal phase at the selected height
may cost only $r$ and must not be assigned a universal $1/2$
separation. Changing ceiling heights between $y$ planes creates
violated $y$ bonds; the same positive comparison excludes that case.

Along factorial sizes all minimizing periods $h+1$ divide $L$.
Their clean row energies are exactly $L\phi_*$, also at ties.
For any other fixed height, the same $L\delta-r$ estimate applies;
the finitely many low candidate heights can be included in the
factorial divisibility threshold. Heights beyond them have an
increasing positive bulk excess. This proves a positive uniform
separation from the set of minimizing clean configurations.
\end{proof}

For the quantum readout introduce an additional classical bit $C$.
Set $C=0$ at $Z_1$ and at every noninterior symbol. On a $z$ bond
whose lower endpoint is interior, require its $C$ to equal the
indicator that the upper endpoint is $W$. This is a two-site
diagonal test and uniquely determines $C$ everywhere in a clean
configuration.

The machine runs on the uncut floor with tape $x$ and time $y$.
Use the limit program for a $\Delta^0_2$ set $A$. The thickness
counter starts at the floor, the output is transported through the
film, and the quantum segment is selected onsite by
$Y_1C(b\mathbin{\mathrm{OR}}f)$. Its length is
$\ell(h)=\bitlength(h)$ and its radices, after settling, are
$2^{\beta A(j)}$. Freeze all unused machine, transport and dynamic
fields, with blank output interpreted as zero. This is exactly the
functional film compiler with its far thickness face replaced by
the selected ceiling; it does not truncate the floor computation.

\begin{theorem}[Wetting transfer]\label{thm:wetting}
For any $\Delta^0_2$ set $A$ containing $0,1,2$ and integer
$\beta\geq1$, this rational
affine family has, for every sufficiently small fixed $g>0$,
\begin{equation}\label{eq:wetting-exact}
 \Delta_\infty(g)=
 \min_{h\in\operatorname*{argmin}\phi_g}
 \Delta^+(K[2^{\beta A(j)}:j<\bitlength(h)]),
\end{equation}
along factorial open cubes. The lower and upper gap exponents
are $\beta$ times the lower and upper densities of $A$.
Every nonnegative finite $\mathsf D_3$ value and ordered
$\mathsf S_3\times\mathsf P_3$ pair occurs. For $g\leq0$ close
to zero the gap tends to zero. The bulk energy density is $g$
for $g<0$ and zero for $g>0$, with a kink at zero.
\end{theorem}
\begin{proof}
Initialize the first three output cells permanently to one before
running the approximation program, as in Section~\ref{sec:film}.
For every read length $\ell\geq3$ the first three radices are then
at least two. The free orbit has $N\geq2(8+4+2)=28$, so its gap,
and hence the full counter gap, is at most
$4\pi^2/N^2\leq\pi^2/196<r$. This is below the conservative
classical separation $r$ of Lemma~\ref{lem:wetting-selection}.
Equality with the orbit gap is not required.
At each fixed selected height only finitely many output bits are
read. They settle as $L\to\infty$, so the full-counter identity
and the all-sector proof of Theorem~\ref{thm:film} give
\eqref{eq:wetting-exact}. The settling time is not generally
computable from $g$; it is not the computable classical selection
threshold.

Since $\log(1/g)=2\log h+\log(1/r)+o(1)$ and
$\log_2(1/D(h))=2\beta\ell(h)\rho_{\ell(h)}(A)+O(\log\ell(h))$,
the normalized gap logarithm is
$\beta\rho_{\ell(h_*(g))}(A)+o(1)$. Every sufficiently large
integer height occurs. Adjacent heights at a tie have either the
same read length or lengths differing by one, which gives the
same endpoint limits. Apply Lemma~\ref{lem:relative-density}
with densities $a/\beta,b/\beta$ and
$\beta\geq\max(1,b)$. Proposition~\ref{prop:critical-ceiling}
gives the matching ceiling.

For $g=0$, the lowest clean corner count is one and minimizing
heights satisfy $h+1>L/2$, equivalently $h\geq\lfloor L/2\rfloor$.
For fixed negative $g$ sufficiently
close to zero, the maximal height $L-2$ minimizes once
$L|g|>r$. The same run charging uses
$1+gL-\phi_g(L-2)\geq3/4$ for sufficiently small $|g|$ to
exclude defective competitors. In either case the minimizing
quantum segments have diverging bitlength, and their main-orbit
upper bound is $\pi^2/\ell^2\to0$, independent of the guessed
bits. Filling the cube at negative $g$ gives energy density $g$;
a bounded-thickness film at fixed positive $g$ has only surface
energy and density zero. The nonnegative dynamic terms have
kernels, so they do not change these densities.
\end{proof}

The volume convention is substantive. At the tie $t_7=1/1152$,
formula~\eqref{eq:wetting-row} selects height $8$ at
$L=72n+8$ and height $7$ at $L=72n+9$. Their read lengths differ,
and a legal word can give different full counter gaps. Thus the
all-size open-cube gap limit need not exist at a tie. Neither an
all-size identity there nor a universal classical separation of
$1/2$ is used. The factorial convention removes this residue
obstruction without changing rational coefficients or the
coupling exponent.

At a single fixed coupling the eventual counter is one finite
algebraic matrix and its gap is individually computable. The
failure is uniform in the coupling. To prove that no general
computable scalar-gap settling modulus exists, choose a realized
exponent outside $\mathsf D_2$. A uniform computable inner-gap
modulus would make the scalar gaps uniformly computable and
Proposition~\ref{prop:critical-ceiling} would put this exponent
in $\mathsf D_2$, a contradiction. No injectivity of the map
from radix vectors to their scalar gaps is assumed.
